\section{Pencil collisions and comparison of local types}
\label{sec:pencil-collision-types}
The normalization criterion above separates phenomena which a
set-theoretic discriminant calculation conflates. In the affine
normalization write $g=hq$ and $v=\lambda+hw$, with $h$ monic of
degree two. The kernel test is
\[
 h\mid qr,\qquad wr\bmod h\text{ constant},\qquad\deg r<2.
\]
The following three cases give concrete local distinctions.

If $g$ has distinct roots, the normalization is unramified; a unique
equal value pair gives a smooth point, and two disjoint equal pairs
give an ordinary double crossing. Three equal values give the three
braid factors $\delta_{12}\delta_{13}(\delta_{13}-\delta_{12})$,
not a normal-crossing divisor in a two-dimensional normal space.
These are exact equations because root labelling is etale here.

For $g=t^3(t-1)$ and $v=t^2$, the scalar subdivisor $h=t^2$ is unique.
One has $q=t(t-1)$ and $w=1$. Divisibility forces the constant term
of $r$ to vanish and the second condition forces its linear term to
vanish. The total pencil divisor is smooth, although its fixed-$D$
fibre has primary equation a unit times $c_{11}^3$.

For $g=t^4$ and $v=t^2$, again $h=t^2$ is the unique scalar
subdivisor. Now $q=t^2$ and $w=1$. The constant $r=1$ satisfies the
kernel test. Thus the total pencil divisor is singular and its
normalization is ramified at the unique preimage. This last example
is asserted to have a unique normalization point, not an unproved
analytic cusp normal form. In contrast, the moving contact-hyperplane
scheme is smooth at all these divisor collisions by
Theorem~\ref{thm:moving-hyperplane}. Its fixed fourfold-contact
transverse algebra is $G_4$, with length six and nilpotency index
five, rather than the pencil's sixth power of a prime divisor.
